\documentclass[11pt,a4paper]{article}
\usepackage[margin=2.5cm]{geometry}
\usepackage{graphicx}
\usepackage{hyperref}
\usepackage{listings}
\usepackage{xcolor}
\usepackage{fancyhdr}
\usepackage{longtable}

\hypersetup{colorlinks=true,linkcolor=blue!60!black,urlcolor=blue!60!black}
\lstset{basicstyle=\ttfamily\small,breaklines=true,frame=single,
        backgroundcolor=\color{gray!8},columns=fullflexible}

\title{\textbf{Reproducible Articulatory-to-Acoustic Reproduction Pipeline}\\[4pt]
\large Documentation of the GitHub repository and the Zenodo archive\\[2pt]
\normalsize Companion material of the article ``Assessing a Mathematical Model
of Syllable Production via CTW Alignment with EMA Data''\\
by F. Berthommier (Univ.\ Grenoble Alpes, CNRS, Grenoble INP, GIPSA-lab, France)\\
(Interspeech 2026 submission 2479, resubmitted to arXiv)}
\author{F. Berthommier -- Gipsa-Lab, CNRS, University Grenoble Alpes}
\date{September 2026}

\pagestyle{fancy}\fancyhf{}
\fancyhead[L]{CTW Reproduction Pipeline}
\fancyhead[R]{\thepage}

\begin{document}
\maketitle

\begin{center}
\fbox{\parbox{0.9\textwidth}{\textbf{Finalization banner.} The cross-references
between the three repositories (GitHub, Zenodo, arXiv) are being finalized.
Insert here the Zenodo DOI, the arXiv identifier and the official
Interspeech-style reference once available
(see \texttt{docs/FINALIZATION-CHECKLIST.md}).}}
\end{center}

\section{Overview}
This documentation explains, for non-expert users, how to install and run the
complete pipeline that reproduces every figure of the article and of its
supplement, the corpus of synthesized short syllables, and the MP4 stimulus
videos. The same document serves as the guide to the components of the Zenodo
archive.

The pipeline operates on articulatory trajectories of spoken syllables. An
EMA (electromagnetic articulography) corpus recorded by \textbf{Kasper and
colleagues}~[1] provides the raw material; the data were kindly
transmitted to the author personally by the authors of that publication
(personal communication, 2026). Reference articulatory parameter tracks
are compared, condition~A (natural alignment) versus condition~B (randomized
control), using \textbf{Canonical Time Warping (CTW)}: a Maeda-style
articulatory model is aligned onto the reference tracks, and the agreement is
quantified by per-parameter correlations with paired statistical tests
(Holm-corrected). The same comparison is carried out on formant tracks
(F1--F3). Synthesis and video generation then turn the aligned parameters into
audible and visible stimuli.

\section{Repository layout (GitHub)}
\begin{longtable}{p{4.2cm}p{10.5cm}}
\texttt{data/} & All data. \texttt{mat/} raw EMA recordings; \texttt{wav/},
\texttt{copywav/} audio; \texttt{maeda/} Maeda model parameters;
\texttt{copyparam/} reference articulatory tracks; \texttt{copyformants/}
reference formant tracks; \texttt{modelctw/}, \texttt{modelctwrandom/} CTW
aligned outputs (conditions A and B); \texttt{paramodelctw/},
\texttt{paramodelctwrandom/} aligned parameter tracks.\\
\texttt{src/} & Reusable library: CTW alignment, correlation computation,
statistical tests.\\
\texttt{figures/} & One self-documented script per article/supplement figure.\\
\texttt{synthesis/} & Short-syllable synthesis pipeline.\\
\texttt{video/} & Audio/video stimulus makers (MP4 reproduction).\\
\texttt{docs/} & This documentation, script mapping, finalization checklist.\\
\texttt{results/} & Outputs regenerated by the scripts (not tracked).\\
\end{longtable}

\section{Installation}
Python~$\geq$~3.9 is required; \texttt{ffmpeg} must be on the \texttt{PATH}
for video reproduction.
\begin{lstlisting}
git clone <REPOSITORY-URL>        % TODO-FINALIZATION: insert final URL
cd <repository>
python -m venv .venv
.venv\Scripts\activate            % Windows  (Linux/macOS: source .venv/bin/activate)
pip install -r requirements.txt
\end{lstlisting}

\section{Reproducing the figures}
Each figure of the article and of the supplement is generated by one script
(see \texttt{docs/mapping.md} for the correspondence with the original
research scripts). Every script saves a PNG in \texttt{results/figures/} and
prints the associated statistics.
\begin{lstlisting}
python figures/fig2a_articulatory_correlations.py    % Figure 2a (article)
python figures/fig2b_formant_correlations.py         % Figure 2b (article)
python figures/fig3_formant_trajectories.py          % Figure 3  (article)
python figures/figs1a_ctw_paths_condition_a.py       % Fig. S1a  (supplement)
python figures/figs1b_ctw_paths_condition_b.py       % Fig. S1b  (supplement)
python figures/figs3_<name>.py                       % Fig. S3   (supplement)
\end{lstlisting}
Each script reads only the \texttt{data/} directory; no network access is
required. All printed output (correlation tables, paired $t$-tests with Holm
correction) is part of the reproducible record and should match the numbers
reported in the article and supplement.

\section{Reproducing the syllable synthesis}
\begin{lstlisting}
python synthesis/run_synthesis.py --condition A --output-dir results/synthesis
\end{lstlisting}
The synthesis pipeline regenerates the short-syllable corpus from the aligned
articulatory parameters. Stimuli are named \texttt{Y0X} with $X \in
[10,16]$ and $Y \in [2,11]$ (63 syllables per condition).

\section{Reproducing the MP4 videos}
The article videos (\texttt{CondA.mp4}, \texttt{CondB.mp4},
\texttt{mixAB.mp4}) are reproduced as follows (see \texttt{video/README.md}
for details):
\begin{lstlisting}
python video/vlam_audio_maker_v2.py     % assemble audio stimuli
python video/vlam_video_maker.py        % render trajectory animations, encode MP4
\end{lstlisting}
Outputs are written to \texttt{results/videos/}. Rendering time depends on
the machine; the codec and frame rate are fixed by the scripts so that the
reproduced videos match the archived ones in content and duration.

\section{Guide to the Zenodo archive}
The Zenodo deposit mirrors and extends the GitHub repository. It contains:
\begin{enumerate}
  \item \textbf{Supplement} --- the (slightly amended) supplement of the
        article, updated so that every described experiment is reproducible
        with the scripts of this repository.
  \item \textbf{Figures} --- all figures of the article and of the supplement
        (PNG), exactly as generated by the \texttt{figures/} scripts.
  \item \textbf{Programs} --- a versioned snapshot of the GitHub repository
        (code without the heavy \texttt{data/} directory).
  \item \textbf{Videos} --- the stimulus videos \texttt{CondA.mp4},
        \texttt{CondB.mp4}, \texttt{mixAB.mp4}.
  \item \textbf{Simulations} --- the intermediate and final outputs of the
        CTW alignment and the statistical analyses (parameter tracks, formant
        tracks, correlation matrices).
\end{enumerate}
Each component is documented in \texttt{GUIDE.md} at the root of the Zenodo
package. Licenses: code, figures, documentation and videos under MIT
(Berthommier, Gipsa-Lab); EMA data under MIT third party (Kasper et al.).

\section{Licenses and citation}
See \texttt{LICENSE}, \texttt{LICENSES/LICENSE-EMA-data.md} and
\texttt{CITATION.cff}. \textbf{TODO-FINALIZATION:} insert the Zenodo DOI, the
arXiv identifier and the official Interspeech-style reference in
\texttt{CITATION.cff}, \texttt{README.md} and this document once the three
deposits are finalized.

\section*{References}
\begin{enumerate}
\item J.~Kasper, C.~A.~Kell, P.~Perrier,
``An informational account of anticipatory coarticulation -- theoretical
considerations and empirical plausibility'',
\emph{Journal of Neurolinguistics}, vol.~78, p.~101298, 2026.
\texttt{doi:10.1016/j.jneuroling.2025.101298}.
The EMA data used in this work were kindly transmitted personally by the
authors of this publication.
\end{enumerate}

\end{document}
